
\medskip

Pour chacune des affirmations suivantes, dire si elle est vraie ou fausse \textbf{en justifiant soigneusement la réponse.}

\medskip

\begin{enumerate}
\item Scratch souhaite rejoindre un ami, mais il a oublié la fin du trajet. Il décide de finir son trajet en prenant, aux intersections, à droite ou à gauche au hasard.

\begin{center}
\psset{unit=1cm}
\begin{pspicture}(0,1)(10,5)
%\psgrid
\psline(0,3.1)(3,3.1)(9,4.6)
\psline(0,2.9)(3,2.9)(9,1.4)
\pspolygon(3.3,3)(6,3.7)(9,3)(6,2.3)
\psline(9,4.4)(6.4,3.8)(9,3.18)
\psline(9,2.82)(6.4,2.2)(9,1.6)
\uput[r](9,4.5){\Large A}
\uput[r](9,3){\Large B}
\uput[r](9,1.5){\Large C}
\rput(0.4,3.25){\includegraphics[width=0.5cm]{chat_Scratch}}
\end{pspicture}
\end{center}

\textbf{Affirmation 1 :} La probabilité qu'il arrive en A, en B ou en C est la même.

\item On suppose qu'une éolienne produit $5$~GWh d'électricité par an et qu'une personne
a besoin de \np{7000}~kWh d'électricité par an. (Wh : Watt-heure)

\textbf{Affirmation 2 :} Une éolienne ne couvre pas les besoins en électricité de \np{1000}~personnes pour un an.

\item Voici quatre nombres : $45\,\%\:;\quad\dfrac{305}{612} \:;\quad 0,5 \:;\quad 730 \times  10^{-3}$.

\textbf{Affirmation 3 :} Ces quatre nombres sont rangés dans l'ordre croissant.

\item L'histogramme ci-dessous représente la répartition des salaires dans une entreprise :

\begin{center}
\psset{xunit=0.01cm,yunit=1cm}
\begin{pspicture}(0,-0.6)(800,5)
\psframe(100,2)\psframe(100,0)(200,5)\psframe(200,0)(300,4)\psframe(300,0)(400,3)
\psframe(400,0)(800,1)
\psframe(500,4)(600,5)\uput[r](600,4.5){\footnotesize 1 bloc = 20 employés}
\uput[d](0,0){\np{1400}}\uput[d](100,0){\np{1500}}\uput[d](200,0){\np{1600}}
\uput[d](300,0){\np{1700}}\uput[d](400,0){\np{1800}}\uput[d](800,0){\np{2200}}
\uput[d](550,-0.4){\footnotesize \emph{source : eduscol.education.fr/ressources-2016}}
\end{pspicture}
\end{center}

\textbf{Affirmation 4 :} Plus de 40\,\% des employés ont un salaire au moins égal à \np{1700}~\euro.
\end{enumerate}

\bigskip

